% R038：纵向残差与最短垂距严格区分；
% 在 x=1cm、y=1.18cm 的画布尺度下，灰线与回归线严格垂直。

\begin{center}
	\begin{tikzpicture}[
			x=1cm,
			y=1.18cm,
			>=Stealth,
			every node/.style={font=\small}
		]

		\path[use as bounding box] (0,0) rectangle (12.4,5.8);

		\node[font=\large\bfseries,text=CExplanation] at (6.2,5.46)
		{最小二乘最小化纵向残差平方，不是点到直线的最短距离};

		% -------------------- 左侧：两种距离 --------------------
		\draw[rounded corners=5pt,fill=blue!2,draw=CExplanation]
		(.15,.28) rectangle (6.02,4.98);

		\draw[->,black!70] (.72,.82)--(5.56,.82) node[right] {$x$};
		\draw[->,black!70] (.72,.82)--(.72,4.35) node[above] {$y$};

		% 回归线
		\draw[CStatement,very thick] (.92,1.25)--(5.28,3.90);

		% 数据点
		\fill[blue!80!black] (3.65,3.88) circle (3.1pt)
		node[above=4pt] {数据点 $(x_i,y_i)$};

		% 红线：固定 x_i 后的纵向残差
		\draw[red!80!black,very thick]
		(3.65,3.88)--(3.65,2.91);

		\node[
			anchor=east,
			text=red!80!black,
			font=\scriptsize
		] at (3.50,3.36)
		{$e_i=y_i-\hat y_i$};

		% 灰线：点到回归线的最短垂距
		\coordinate (foot) at (4.1925,3.2390);

		\draw[gray!80,very thick]
		(3.65,3.88)--(foot);

		\node[
			anchor=west,
			text=gray!80,
			font=\scriptsize
		] at (4.08,2.90)
		{灰：最短垂距};

		% 直角标记
		\draw[gray!80,thick]
		(4.10,3.35)--(3.97,3.27)--(4.06,3.16);

		\node[text=CExplanation,font=\bfseries] at (3.10,.54)
		{红：固定 $x_i$ 的纵向差};

		% -------------------- 右侧：为什么平方 --------------------
		\draw[rounded corners=5pt,fill=orange!3,draw=orange!80!black]
		(6.40,.28) rectangle (12.25,4.98);

		\node[font=\bfseries,text=orange!85!black] at (9.32,4.48)
		{为什么要平方？};

		% 正负误差分开显示
		\draw[red!75!black,->,thick]
		(7.90,3.64)--(9.05,3.64)
		node[midway,above=2pt] {$+2$};

		\draw[red!75!black,->,thick]
		(10.74,3.64)--(9.59,3.64)
		node[midway,above=2pt] {$-2$};

		\node at (9.32,3.08)
		{$+2+(-2)=0$：会抵消};

		% 三个正方形：边长按 1:2:3，面积按 1:4:9
		% y 单位为 1.18cm，因此纵坐标高度除以 1.18

		\draw[fill=orange!25,draw=orange!85!black]
		(7.15,1.42) rectangle (7.65,1.844);
		\node at (7.40,1.632) {$1$};

		\draw[fill=orange!25,draw=orange!85!black]
		(8.82,1.42) rectangle (9.82,2.267);
		\node at (9.32,1.844) {$4$};

		\draw[fill=orange!25,draw=orange!85!black]
		(10.40,1.42) rectangle (11.90,2.691);
		\node at (11.15,2.056) {$9$};

		\node at (7.40,1.07) {$1^2=1$};
		\node at (9.32,1.07) {$2^2=4$};
		\node at (11.15,1.07) {$3^2=9$};

		\node[font=\small\bfseries] at (9.32,.58)
		{不抵消，且大误差惩罚更重};

	\end{tikzpicture}
\end{center}